\chapter{Словарь}

\begin{description}[style=nextline, leftmargin=2em, itemsep=5pt]
  \item[ADB] Android Debug Bridge. Отладочный канал по USB, даёт оболочку и копирование файлов. Запущен и на дроне, и на модуле.
  \item[AprilTag] Семейство плоских меток-квадратов для компьютерного зрения. Наши ковры используют \texttt{tag36h11}.
  \item[ArUco] Другое семейство меток из OpenCV. Штатный модуль его не распознаёт.
  \item[awlink] Протокол кадров между дроном и приложениями (глава~\ref{ch:app}).
  \item[CRC] Контрольная сумма для проверки целостности кадра.
  \item[Зуммер] Пищалка с четырьмя проводами на разъёме U0TX/U0RX модуля; понимает кадр \texttt{AD FA 01 нота длительность сумма} (глава~\ref{ch:pult}).
  \item[Ретранслятор] Наша маленькая программа на дроне: принимает UDP-пакеты с компьютера и пересылает их на UART модуля (глава~\ref{ch:pult}).
  \item[Пульт (USB)] Пульт Wing с USB-кабелем (CP210x, COM3): с компьютера через него можно читать состояние дрона и управлять им.
  \item[imlib] Библиотека компьютерного зрения из проекта OpenMV; на ней построен детектор модуля.
  \item[Melis] Операционная система реального времени на ядре RISC-V дрона C10.
  \item[MCRF4XX] Вариант CRC-16 с начальным значением \texttt{FFFF}, которым защищены кадры awlink.
  \item[\texttt{pilot}, \texttt{pilot\_f}] Программы-автопилоты на дроне Wing (\texttt{\_f}: вариант, запущенный на нашем дроне).
  \item[ToF] Time of Flight. Датчик расстояния по времени пролёта света; в модуле VI5300.
  \item[UART] Простой последовательный интерфейс. На дроне соединяет автопилот с модулем (\texttt{ttyS0}, 115\,200 бод).
  \item[V831] Процессор Allwinner с нейроускорителем, на котором работает позиционный модуль.
  \item[WS2812] Адресуемые светодиоды с порядком цветов зелёный, красный, синий.
  \item[Формация] Режим нескольких дронов, в которых дрон подключается к сети как клиент.
\end{description}
